\section{Lógica proposicional}\label{sec:logicaProposicional}\index{lógica!proposicional}
Nosso objetivo nesta seção é fixar três ideias básicas: a definição de fórmula proposicional, a introdução da noção de dedução formal, e, por fim, como essa noção sintática de dedução se relaciona com a noção semântica de tautologia, que também será introduzida nesta seção.

\subsection{Fórmulas proposicionais e convenções de escrita}\label{subsec:formulasProposicionaisConvencoes}

Iniciaremos definindo a linguagem da lógica proposicional e suas fórmulas.

\begin{metadefinition}[Fórmulas proposicionais]\label{mdef:formulasProposicionais}\index{fórmula!proposicional}
    O léxico da lógica proposicional é o léxico para a notação polonesa cujos símbolos básicos são:
    \begin{itemize}
        \item $\neg$ (negação, unário);
        \item $\wedge$ (conjunção, binário);
        \item $\vee$ (disjunção, binário);
        \item $\rightarrow$ (condicional, binário);
        \item $\leftrightarrow$ (bicondicional, binário);
        \item uma infinidade (enumerável) de símbolos de \emph{variáveis proposicionais} $P_0, P_1, \dots$ (símbolos $0$-ários).
    \end{itemize}
    As \emph{fórmulas proposicionais} são as expressões bem formadas nesse léxico.
\end{metadefinition}

As variáveis proposicionais são símbolos que podem ser interpretados como proposições, ou seja, como afirmações que podem ser verdadeiras ou falsas.

\begin{example}\label{ex:formulasProposicionais}Abaixo, daremos alguns exemplos de fórmulas proposicionais na notação polonesa.
    \begin{itemize}
        \item $\neg P_0$ é uma fórmula proposicional.
        \item $\wedge P_1 P_2$ é uma fórmula proposicional.
        \item $\rightarrow \neg P_3 P_4$ é uma fórmula proposicional.
        \item $\leftrightarrow \neg \wedge P_5 P_6 \neg P_7$ é uma fórmula proposicional.
        \item $\rightarrow \vee P_0 P_7 \wedge P_2 P_1$ é uma fórmula proposicional.
    \end{itemize}
\end{example}
Escrever fórmulas proposicionais na notação polonesa, apesar de formalmente correto, é pouco prático para a maioria dos propósitos, e pode se transformar em um obstáculo didático.
Por isso, utilizaremos as seguintes convenções.
\begin{convention}\label{conv:convencoesEscritaPC}Considere fórmulas proposicionais $\phi$ e $\psi$.
Utilizaremos as seguintes convenções de escrita.
\begin{itemize}
    \item $\phi \wedge \psi$ em vez de ${\wedge}\phi\psi$;
    \item $\phi \vee \psi$ em vez de ${\vee}\phi\psi$;
    \item $\phi \rightarrow \psi$ em vez de ${\rightarrow}\phi\psi$;
    \item $\phi \leftrightarrow \psi$ em vez de ${\leftrightarrow}\phi\psi$.
\end{itemize}
Utilizaremos parênteses para eliminar ambiguidades sempre que for necessário.

Convencionaremos, a fim de eliminar boa parte do uso de parenteses, que a precedência do conectivo $\neg$ é maior do que a de todos os outros, de modo que, por exemplo, $\neg P_0 \vee P_1$ seja interpretado como $(\neg P_0) \vee P_1$.

Adicionalmente, convencionamos que $\rightarrow$ e $\leftrightarrow$ têm precedência menor do que $\wedge$ e $\vee$, de modo que, por exemplo, $P_0 \wedge P_1 \rightarrow P_2$ seja interpretado como $(P_0 \wedge P_1) \rightarrow P_2$.
\end{convention}

Não adotaremos, neste ponto, nenhuma convenção associativa para conectivos binários repetidos.
Assim, expressões como $P_0\wedge P_1\wedge P_2$ deverão receber parênteses sempre que forem utilizadas.
Após demonstrarmos as equivalências associativas pertinentes, poderemos omitir tais parênteses quando isso não gerar ambiguidade.

No exemplo \ref{ex:formulasProposicionais}, as fórmulas proposicionais seriam escritas, respectivamente, como:
\begin{itemize}
    \item $\neg P_0$;
    \item $P_1 \wedge P_2$;
    \item $\neg P_3 \rightarrow P_4$;
    \item $\neg(P_5 \wedge P_6) \leftrightarrow \neg P_7$;
    \item $(P_0 \vee P_7) \rightarrow (P_2 \wedge P_1)$.
\end{itemize}

\subsection{Deduções formais}\label{subsec:deducoesFormaisProposicionais}

Intuitivamente, se temos que $P_0$ é verdade e $P_0\rightarrow P_1$ é verdade, então $P_1$ é verdade.
Tal inferência é chamada de \emph{modus ponens}.
Além disso, intuitivamente, $P_3\land P_4\rightarrow P_3$ é uma fórmula verdadeira, independentemente do que sejam $P_3$ e $P_4$.

Para dar sentido a estes fatos intuitivos, introduziremos a noção de dedução formal da lógica proposicional.
De forma geral, dedução formal é uma sequência de fórmulas válidas de modo que cada fórmula da sequência ou é um axioma do sistema ou é obtida a partir de uma regra de inferência aplicada a fórmulas anteriores.

Iniciaremos delimitando axiomas para a lógica proposicional clássica.
Não trataremos, neste texto, de outros sistemas de lógica proposicional, como a lógica proposicional intuicionista ou a lógica proposicional mínima.

\begin{metadefinition}[Axiomas da lógica proposicional]\label{mdef:axiomasLogicaProposicional}\index{axioma!lógica proposicional}
    Os axiomas da Lógica Proposicional Clássica são as fórmulas que podem ser escritas de alguma das seguintes formas, em que $\phi$, $\psi$ e $\chi$ são outras fórmulas proposicionais.

    Mais precisamente, cada item abaixo é um \emph{esquema de axiomas}.
    Sempre que substituímos $\phi$, $\psi$ e $\chi$ por fórmulas proposicionais quaisquer, obtemos um axioma da lógica proposicional.

    \begin{enumerate}[label=(PC\arabic*)]
        \item $\phi \rightarrow (\psi \rightarrow \phi)$;
        \item $(\phi \rightarrow (\psi \rightarrow \chi)) \rightarrow ((\phi \rightarrow \psi) \rightarrow (\phi \rightarrow \chi))$;
        \item $(\phi \rightarrow \psi) \rightarrow ((\phi \rightarrow \neg \psi) \rightarrow \neg \phi)$;
        \item $(\phi \land \psi) \rightarrow \phi$;
        \item $(\phi \land \psi) \rightarrow \psi$;
        \item $\phi \rightarrow (\psi \rightarrow (\phi \land \psi))$;
        \item $\phi \rightarrow (\phi \lor \psi)$;
        \item $\psi \rightarrow (\phi \lor \psi)$;
        \item $(\phi \rightarrow \chi) \rightarrow ((\psi \rightarrow \chi) \rightarrow ((\phi \lor \psi) \rightarrow \chi))$;
        \item $(\psi \leftrightarrow \phi) \rightarrow (\psi \rightarrow \phi)$;
        \item $(\psi \leftrightarrow \phi) \rightarrow (\phi \rightarrow \psi)$.
        \item $(\psi \rightarrow \phi) \rightarrow ((\phi \rightarrow \psi) \leftrightarrow (\psi \leftrightarrow \phi))$.
        \item $\neg \neg \phi \rightarrow \phi$.
    \end{enumerate}
\end{metadefinition}


Agora estamos prontos para definir o que é uma dedução formal na lógica proposicional.

\begin{metadefinition}[Deduções formais na lógica proposicional]\label{mdef:deducoesFormaisLogicaProposicional}\index{dedução formal!lógica proposicional}
    Considere uma coleção de fórmulas proposicionais $\Gamma$.

    Uma \emph{dedução formal na lógica proposicional a partir de $\Gamma$} é uma sequência de fórmulas proposicionais $\phi_0, \dots, \phi_n$ tal que, para cada $i \leq n$:
    \begin{itemize}
        \item $\phi_i$ é uma fórmula de $\Gamma$, ou
        \item $\phi_i$ é um axioma da lógica proposicional, ou
        \item existem $j, k < i$ tais que $\phi_j$ é da forma $\phi_k \rightarrow \phi_i$.
    \end{itemize}
    
    Dizemos que uma fórmula \emph{$\psi$ é deduzida a partir de $\Gamma$}, ou que \emph{$\psi$ é uma consequência lógica de $\Gamma$}, ou que \emph{$\psi$ é um teorema de $\Gamma$}, e escrevemos $\Gamma \vdash_P \psi$, se existe uma dedução formal $\phi_0, \dots, \phi_n$ a partir de $\Gamma$ tal que $\phi_n$ é $\psi$.

    Na notação do parágrafo anterior, $\Gamma$ é a coleção vazia de fórmulas proposicionais, então dizemos que \emph{$\psi$ é um teorema da lógica proposicional}, e escrevemos $\vdash_P \psi$.
\end{metadefinition}

Vamos começar com dois exemplos simples.

\begin{example}Se $\Gamma$ é uma coleção de fórmulas proposicionais e $\phi$ é uma fórmula de $\Gamma$, então $\Gamma \vdash_P \phi$.
\end{example}

\begin{example}$\vdash_P (P_0 \rightarrow (P_1 \rightarrow P_0))$, pois $P_0 \rightarrow (P_1 \rightarrow P_0)$ é um axioma da lógica proposicional.
\end{example}

Também daremos a seguinte observação simples.

\begin{metaproposition}\label{metaprop:PCmonotonicidade}Sejam $\Gamma$ e $\Gamma'$ coleções de fórmulas proposicionais tais que $\Gamma \subseteq \Gamma'$ (ou seja, tais que toda fórmula de $\Gamma$ é uma fórmula de $\Gamma'$).
Se $\Gamma \vdash_P \phi$, então $\Gamma' \vdash_P \phi$.
\end{metaproposition}
\begin{proof}Se $\Gamma \vdash_P \phi$, então existe uma dedução formal $\phi_0, \dots, \phi_n$ a partir de $\Gamma$ tal que $\phi_n$ é $\phi$.
É imediato da definição de dedução formal que $\phi_0, \dots, \phi_n$ é uma dedução formal a partir de $\Gamma'$, assim, $\Gamma' \vdash_P \phi$.
\end{proof}

Logo, parece ser relevante saber quais as fórmulas proposicionais $\phi$ tais que $\vdash_P \phi$, uma vez que elas são consequências lógicas de qualquer coleção de fórmulas proposicionais.

Iniciaremos com o seguinte.

\begin{metaproposition}\label{metaprop:PCreflexividade}Se $\phi$ é qualquer fórmula proposicional, então $\vdash_P \phi \rightarrow \phi$.
\end{metaproposition}
\begin{proof}Vamos construir uma dedução formal de $\phi \rightarrow \phi$.
    Considere:

    \begin{enumerate}
        \item $(\phi \rightarrow ((\phi \rightarrow \phi) \rightarrow \phi)) \rightarrow ((\phi \rightarrow (\phi \rightarrow \phi)) \rightarrow (\phi \rightarrow \phi))$ (PC2 com $\phi$, $\phi\rightarrow \phi$ e $\phi$);
        \item $\phi \rightarrow ((\phi \rightarrow \phi) \rightarrow \phi)$ (PC1 com $\phi$ e $\phi\rightarrow \phi$);
        \item $(\phi \rightarrow (\phi \rightarrow \phi)) \rightarrow (\phi \rightarrow \phi)$ (MP 1 e 2);
        \item $\phi \rightarrow (\phi \rightarrow \phi)$ (PC1 com $\phi$ e $\phi$);
        \item $\phi \rightarrow \phi$ (MP 3 e 4).
    \end{enumerate}
\end{proof}
Apesar de curta, a dedução formal da metaproposição anterior é um pouco difícil de ser encontrada.
Seria muito útil ter métodos mais simples para demonstrar que uma fórmula proposicional é um teorema da lógica proposicional.
O primeiro método, simples, é uma abstração do modus ponens.
\begin{metaproposition}\label{metaprop:PCmodusPonens}Considere fórmulas proposicionais $\phi$ e $\psi$, e uma coleção de fórmulas proposicionais $\Gamma$.
Se $\Gamma \vdash_P \phi$ e $\Gamma \vdash_P \phi \rightarrow \psi$, então $\Gamma \vdash_P \psi$.
\end{metaproposition}
\begin{proof}Como $\Gamma \vdash_P \phi$, então existe uma dedução formal $\phi_0, \dots, \phi_n$ a partir de $\Gamma$ tal que $\phi_n$ é $\phi$.
Como $\Gamma \vdash_P \phi \rightarrow \psi$, então existe uma dedução formal $\psi_0, \dots, \psi_m$ a partir de $\Gamma$ tal que $\psi_m$ é $\phi \rightarrow \psi$.
Concatenando as deduções formais $\phi_0, \dots, \phi_n$ e $\psi_0, \dots, \psi_m$, e adicionando $\psi$ ao final da sequência, obtemos uma dedução formal a partir de $\Gamma$ tal que $\psi$ é a última fórmula da dedução formal, ou seja, $\Gamma \vdash_P \psi$.
\end{proof}

Um método central para trabalhar com deduções lógicas é o chamado Teorema da Dedução.

Primeiro, uma convenção.
\begin{convention}Considere uma coleção de fórmulas proposicionais $\Gamma$ e fórmulas proposicionais $\phi_1, \dots, \phi_n$ e $\psi$.
Então $\Gamma\cup \{\phi_1, \dots, \phi_n\}$ denota a coleção de fórmulas proposicionais que contém todas as fórmulas de $\Gamma$ e as fórmulas $\phi_1, \dots, \phi_n$.
\end{convention}

\begin{metatheorem}[Teorema da Dedução]\label{thm:teoremaDeducao}\index{teorema!da dedução}
    Considere uma coleção de fórmulas proposicionais $\Gamma$ e fórmulas proposicionais $\phi$ e $\psi$.
    Então $\Gamma \cup \{\phi\} \vdash_P \psi$ se, e somente se, $\Gamma \vdash_P \phi \rightarrow \psi$.
\end{metatheorem}
\begin{proof}
    Suponha primeiramente que $\Gamma \vdash_P \phi \rightarrow \psi$.
    Então, pela Metaproposição~\ref{metaprop:PCmonotonicidade}, $\Gamma \cup \{\phi\} \vdash_P \phi \rightarrow \psi$.
    Como $\phi$ está em $\Gamma \cup \{\phi\}$, então $\Gamma \cup \{\phi\} \vdash_P \phi$.
    Assim, por modus ponens, $\Gamma \cup \{\phi\} \vdash_P \psi$.

    Agora suponha que $\Gamma \cup \{\phi\} \vdash_P \psi$.
    Então existe uma dedução formal $\psi_0, \dots, \psi_n$ a partir de $\Gamma \cup \{\phi\}$ tal que $\psi_n$ é $\psi$.
    Por indução para $i\leq n$, veremos que $\Gamma \vdash_P \phi \rightarrow \psi_i$ para cada $i\leq n$.

    Suponha que a hipótese indutiva é verdadeira para cada $j < i$ e que $i\leq n$.

    \textbf{Caso 1:} $\psi_i$ é uma fórmula de $\Gamma$ ou um axioma da lógica proposicional.
    Então $\Gamma \vdash_P \psi_i$.
    Temos que $\Gamma \vdash_P \psi_i \rightarrow (\phi \rightarrow \psi_i)$, pois $\psi_i \rightarrow (\phi \rightarrow \psi_i)$ é um axioma da lógica proposicional.
    Assim, por modus ponens, $\Gamma \vdash_P \phi \rightarrow \psi_i$.

    \textbf{Caso 2:} $\psi_i$ é $\phi$.
    Então $\Gamma \vdash_P \phi \rightarrow \phi$ pela Metaproposição~\ref{metaprop:PCreflexividade}.

    \textbf{Caso 3:} existem $j, k < i$ tais que $\psi_j$ é da forma $\psi_k \rightarrow \psi_i$.
    Então, pela hipótese indutiva, $\Gamma \vdash_P \phi \rightarrow \psi_j$ e $\Gamma \vdash_P \phi \rightarrow \psi_k$.
    Temos que $\Gamma \vdash_P (\phi \rightarrow \psi_j) \rightarrow ((\phi \rightarrow \psi_k) \rightarrow (\phi \rightarrow \psi_i))$, pois esta fórmula é um axioma da lógica proposicional (PC2).
    Assim, por modus ponens, $\Gamma \vdash_P (\phi \rightarrow \psi_k) \rightarrow (\phi \rightarrow \psi_i)$.
    Aplicando modus ponens novamente, $\Gamma \vdash_P \phi \rightarrow \psi_i$.

    Portanto, por indução, $\Gamma\vdash_P \phi\rightarrow \psi_i$ para todo $i\leq n$.
    Em particular, como $\psi_n$ é $\psi$, temos $\Gamma\vdash_P \phi\rightarrow\psi$.
\end{proof}

Como aplicação, provaremos a recíproca de PC13.

\begin{metaproposition}\label{metaprop:PC13reciproca}Seja $\phi$ uma fórmula proposicional.
    Então $\vdash_P \phi \rightarrow \neg\neg \phi$.
\end{metaproposition}
\begin{proof}
    É suficiente ver que $\{\phi\} \vdash_P \neg\neg \phi$.

    Primeiro, $\{\phi\} \vdash_P \phi\rightarrow (\neg \phi \rightarrow \phi)$, pois $\phi\rightarrow (\neg \phi \rightarrow \phi)$ é um axioma da lógica proposicional (PC1).
    Por modus ponens, segue que:
    \[
        \{\phi\} \vdash_P \neg \phi \rightarrow \phi.
    \]

    Temos também que:
    \[
        \{\phi\} \vdash_P \neg\phi\rightarrow \neg\phi
    \]
    Por PC3, temos que:
    \[
        \{\phi\} \vdash_P (\neg \phi \rightarrow \phi) \rightarrow ((\neg \phi \rightarrow \neg \phi) \rightarrow \neg \neg \phi).
    \]

    Aplicando modus ponens duas vezes, segue que $\{\phi\} \vdash_P \neg\neg \phi$.

    Pelo Teorema da Dedução, $\vdash_P \phi \rightarrow \neg\neg \phi$.
\end{proof}

Outro teorema relevante cuja demonstração é facilitada pelo Teorema da Dedução é o Princípio da explosão.
Em suma, ele diz que se for derivada uma contradição, então qualquer fórmula proposicional pode ser derivada.

\begin{metaproposition}[Princípio da explosão]\label{metaprop:principioExplosaoProposicional}Sejam $\phi$ e $\psi$ fórmulas proposicionais.
    Então $\vdash_P \phi\rightarrow(\neg \phi \rightarrow \psi)$.
\end{metaproposition}
\begin{proof}
    É suficiente, por duas aplicações do Teorema da Dedução, ver que $\{\phi, \neg \phi\} \vdash_P \psi$.

    Por PC1, segue que $\{\phi, \neg \phi\} \vdash_P \phi\rightarrow(\psi\rightarrow \phi)$ e que $\{\phi, \neg \phi\} \vdash_P \neg \phi \rightarrow (\psi \rightarrow \neg \phi)$.
    Por modus ponens, $\{\phi, \neg \phi\} \vdash_P \psi \rightarrow \phi$ e $\{\phi, \neg \phi\} \vdash_P \psi \rightarrow \neg \phi$.
    Por PC3, temos:

    \[
        \{\phi, \neg \phi\} \vdash_P (\psi \rightarrow \phi) \rightarrow ((\psi \rightarrow \neg \phi) \rightarrow \neg \psi).
    \]
    Aplicando modus ponens duas vezes, segue que $\{\phi, \neg \phi\} \vdash_P \neg \psi$.
    Repetindo o argumento, segue que $\{\phi, \neg \phi\} \vdash_P \neg \neg \psi$.
    Por PC13 e modus ponens, $\{\phi, \neg \phi\} \vdash_P \psi$.
\end{proof}

Na lógica clássica, também temos acesso ao chamado \emph{modus tollens}.
\begin{metaproposition}[Modus tollens]\label{metaprop:modusTollensProposicional}
    Sejam $\phi$ e $\psi$ fórmulas proposicionais e $\Gamma$ uma coleção de fórmulas proposicionais.
    \begin{enumerate}[label=(\alph*)]
        \item $\vdash_P (\phi \rightarrow \psi) \rightarrow ((\neg \psi) \rightarrow (\neg \phi))$.
        \item $\vdash_P (\neg \psi \rightarrow \neg \phi) \rightarrow (\phi \rightarrow \psi)$.
        \item $\Gamma \vdash_P \phi \rightarrow \psi$ se, e somente se, $\Gamma \vdash_P \neg \psi \rightarrow \neg \phi$.
    \end{enumerate}
\end{metaproposition}
\begin{proof}
    (a) É suficiente ver que $\{\phi \rightarrow \psi, \neg \psi\} \vdash_P \neg \phi$.

    Temos que
    \[
        \{\phi \rightarrow \psi, \neg \psi\} \vdash_P
        (\phi \rightarrow \psi) \rightarrow ((\phi \rightarrow \neg \psi) \rightarrow \neg \phi),
    \]
    pelo axioma PC3.

    Temos que $\{\phi \rightarrow \psi, \neg \psi\} \vdash_P\phi \rightarrow \psi$, assim, por modus ponens, $\{\phi \rightarrow \psi, \neg \psi\} \vdash_P (\phi \rightarrow \neg \psi)\rightarrow \neg \phi$.

    Temos que $\{\phi \rightarrow \psi, \neg \psi\} \vdash_P \neg \psi\rightarrow(\phi \rightarrow \neg \psi)$, pelo axioma PC1, e $\{\phi \rightarrow \psi, \neg \psi\} \vdash_P \neg \psi$, assim, por modus ponens, $\{\phi \rightarrow \psi, \neg \psi\} \vdash_P \phi \rightarrow \neg \psi$.

    Logo, novamente por modus ponens, $\{\phi \rightarrow \psi, \neg \psi\} \vdash_P \neg \phi$.

    (b) Pelo Teorema da Dedução, é suficiente ver que $\{\neg \psi \rightarrow \neg \phi, \phi\} \vdash_P \psi$.
    Pelo item (a),
    \[
        \{\neg \psi \rightarrow \neg \phi, \phi\}\vdash_P
        (\neg \psi \rightarrow \neg \phi) \rightarrow (\neg\neg \phi \rightarrow \neg\neg \psi).
    \]
    Por modus ponens,
    \[
        \{\neg \psi \rightarrow \neg \phi, \phi\}\vdash_P
        \neg\neg \phi \rightarrow \neg\neg \psi.
    \]
    Pela Metaproposição~\ref{metaprop:PC13reciproca} e modus ponens,
    \[
        \{\neg \psi \rightarrow \neg \phi, \phi\}\vdash_P \neg\neg\phi.
    \]
    Logo,
    \[
        \{\neg \psi \rightarrow \neg \phi, \phi\}\vdash_P \neg\neg\psi.
    \]
    Por PC13 e modus ponens,
    \[
        \{\neg \psi \rightarrow \neg \phi, \phi\}\vdash_P \psi.
    \]

    (c) A implicação da esquerda para a direita segue do item (a), da monotonicidade e de modus ponens.
    A implicação da direita para a esquerda segue do item (b), da monotonicidade e de modus ponens.
\end{proof}
\subsection{Tautologias}\label{subsec:tautologias}
Até aqui, tratamos apenas de uma noção sintática: uma fórmula é dedutível quando aparece ao final de uma sequência construída segundo regras formais.
Agora introduziremos uma noção semântica, baseada em valores de verdade.
O objetivo será mostrar que, na lógica proposicional clássica, essas duas perspectivas coincidem.

\begin{metadefinition}[Valoração]\label{mdef:valoracao}\index{valoração}\index{tabela-verdade}
    Considere uma coleção finita de variáveis proposicionais, $\mathcal P$.
    Uma \emph{valoração} de variáveis proposicionais é um mapeamento $v$ que atribui a cada variável proposicional de $\mathcal P$ um valor de verdade, ou seja, $T$ (verdadeiro) ou $F$ (falso).

    Seja $\phi$ uma fórmula proposicional.
    Dizemos que \emph{$v$ é uma valoração de $\phi$} se $v$ é uma valoração que atribui um valor de verdade a cada variável proposicional que ocorre em $\phi$.
\end{metadefinition}

$T$ e $F$ são símbolos arbitrários que não fazem parte da linguagem da lógica de primeira ordem fixados para intuitivamente denotar os valores de verdade verdadeiro e falso, respectivamente.

Notemos que se $v$ é uma valoração de $\phi$, então $v$ é uma valoração de cada subfórmula de $\phi$, de modo que a seguinte definição faz sentido.

A seguir, expandiremos o valor de verdade de uma fórmula proposicional $\phi$ sob uma valoração $v$ a partir do valor de verdade que $v$ atribui a cada variável proposicional que ocorre em $\phi$.

Para tanto, utilizaremos as seguintes tabelas.

\begin{metadefinition}
    Definimos as seguintes operações lógicas que agem nos símbolos $T$ e $F$.
    Nas tabelas abaixo, $A$ e $B$ denotam valores de verdade, ou seja, $T$ ou $F$.
    Cada linha corresponde a um caso possível de atribuição de valores de verdade a $A$ e $B$, e a última coluna corresponde ao valor de verdade da operação lógica aplicada a $A$ e $B$.
    \begin{itemize}
       \item A operação unária $\neg$ é definida por:
        \begin{center}
            \begin{tabular}{c|c}
                $A$ & $\neg A$ \\
                \hline
                $T$ & $F$ \\
                $F$ & $T$
            \end{tabular}
        \end{center}
        \item A operação binária $\wedge$ é definida por:
        \begin{center}
            \begin{tabular}{c|c|c}
                $A$ & $B$ & $A\wedge B$ \\
                \hline
                $T$ & $T$ & $T$ \\
                $T$ & $F$ & $F$ \\
                $F$ & $T$ & $F$ \\
                $F$ & $F$ & $F$
            \end{tabular}
        \end{center}
        \item A operação binária $\vee$ é definida por:
        \begin{center}
            \begin{tabular}{c|c|c}
                $A$ & $B$ & $A\vee B$ \\
                \hline
                $T$ & $T$ & $T$ \\
                $T$ & $F$ & $T$ \\
                $F$ & $T$ & $T$ \\
                $F$ & $F$ & $F$
            \end{tabular}
        \end{center}
        \item A operação binária $\rightarrow$ é definida por:
        \begin{center}
            \begin{tabular}{c|c|c}
                $A$ & $B$ & $A\rightarrow B$ \\
                \hline
                $T$ & $T$ & $T$ \\
                $T$ & $F$ & $F$ \\
                $F$ & $T$ & $T$ \\
                $F$ & $F$ & $T$
            \end{tabular}
        \end{center}
        \item A operação binária $\leftrightarrow$ é definida por:
        \begin{center}
            \begin{tabular}{c|c|c}
                $A$ & $B$ & $A\leftrightarrow B$ \\
                \hline
                $T$ & $T$ & $T$ \\
                $T$ & $F$ & $F$ \\
                $F$ & $T$ & $F$ \\
                $F$ & $F$ & $T$
            \end{tabular}
        \end{center}
    \end{itemize}
\end{metadefinition}

\begin{metadefinition}\label{mdef:valorVerdadeFormulaProposicional}Seja $v$ uma valoração de variáveis proposicionais.
    Além dos símbolos $T$ e $F$ fixados acima, fixe também um símbolo $I$ (que abrevia \emph{indefinido}).

    Recursivamente na complexidade de $\phi$, definimos o valor de verdade de $\phi$ sob $v$, denotado por $v(\phi)$, a partir do seguinte algoritmo.
    Se $v$ não é uma valoração para $\phi$, então $v(\phi)$ é $I$.
    Caso contrário, procedemos como segue.
    \begin{enumerate}
        \item Se $\phi$ é uma variável proposicional $P$, então $v(\phi)$ é o valor de verdade que $v$ atribui a $P$.
        \item Se $\phi$ é $\neg \psi$, então $v(\phi)$ é $\neg v(\psi)$.
        \item Se $\phi$ é $ \psi \square \eta$, em que $\square$ é $\vee$, $\wedge$, $\rightarrow$ ou $\leftrightarrow$, então $v(\phi)$ é $v(\psi) \square v(\eta)$.
    \end{enumerate}
    Dizemos que $\phi$ é uma \emph{tautologia} se $v(\phi)$ é $T$ para toda valoração $v$ de $\phi$.
\end{metadefinition}

Uma primeira observação é que valorações para $\phi$ que coincidem em todas as variáveis proposicionais que ocorrem em $\phi$ atribuem o mesmo valor de verdade a $\phi$.
Deste modo, valores eventualmente atribuídos a variáveis proposicionais que não ocorrem em $\phi$ não afetam o valor de verdade de $\phi$.

\begin{metalemma}
    Sejam $v$ e $v'$ valorações de variáveis proposicionais tais que, para cada variável proposicional $P$ que ocorre em $\phi$, $v(P)$ é igual a $v'(P)$.
    Então $v(\phi)$ é igual a $v'(\phi)$, e não é $I$.
\end{metalemma}
\begin{proof}
    Agimos por indução na complexidade de $\phi$.

    \textbf{Caso 1:} $\phi$ é uma variável proposicional $P$.
    Então $v(\phi)$ é o valor de verdade que $v$ atribui a $P$, e $v'(\phi)$ é o valor de verdade que $v'$ atribui a $P$.
    Por hipótese, $v(P)$ é igual a $v'(P)$, e não é $I$.

    \textbf{Caso 2:} $\phi$ é $\neg \psi$.
    Por hipótese indutiva, $v(\psi)$ é igual a $v'(\psi)$, e não é $I$.
    Logo, $v(\phi)$ é $\neg v(\psi)$, e $v'(\phi)$ é $\neg v'(\psi)$, de modo que $v(\phi)$ é igual a $v'(\phi)$.

    \textbf{Caso 3:} $\phi$ é $ \psi \square \eta$, em que $\square$ é $\vee$, $\wedge$, $\rightarrow$ ou $\leftrightarrow$.
    Por hipótese indutiva, $v(\psi)$ é igual a $v'(\psi)$, e não é $I$, e $v(\eta)$ é igual a $v'(\eta)$, e não é $I$.
    Logo, $v(\phi)$ é $v(\psi) \square v(\eta)$, e $v'(\phi)$ é $v'(\psi) \square v'(\eta)$, de modo que $v(\phi)$ é igual a $v'(\phi)$.
\end{proof}

Veremos inicialmente que todo teorema da lógica é uma tautologia.
Para tanto, precisaremos de mais algumas definições.

\begin{metadefinition}Seja $\mathcal P$ uma coleção finita de variáveis proposicionais.
    Uma atribuição de fórmulas para $\mathcal P$ é um mapeamento $\rho$ que atribui a cada variável proposicional $P$ de $\mathcal P$ uma fórmula proposicional $\rho(P)$.

    Se $\theta$ é uma fórmula proposicional, dizemos que \emph{$\rho$ é uma atribuição para $\theta$} se $\rho$ é uma atribuição de fórmulas para cada variável proposicional que ocorre em $\theta$.
\end{metadefinition}

\begin{metadefinition}[Substituição proposicional]\label{mdef:substituicaoProposicional}\index{substituição!proposicional}
   Considere uma coleção finita de variáveis proposicionais $\mathcal P$ e $\rho$ uma atribuição de fórmulas para $\mathcal P$.

   Recursivamente na complexidade de uma fórmula proposicional $\theta$, definimos a substituição de $\theta$ via $\rho$, denotada por $\subst(\theta, \rho)$, a partir do seguinte algoritmo.
   Se $\rho$ não é uma atribuição para $\theta$, então $\subst(\theta, \rho)$ é $I$.
   Caso contrário, procedemos como nos itens abaixo.
    \begin{itemize}
        \item Se $\theta$ é uma variável proposicional $P$, então $\subst(\theta, \rho)$ é $\rho(P)$.
        \item Se $\theta$ é $\neg \psi$, então $\subst(\theta, \rho)$ é $\neg \subst(\psi, \rho)$.
        \item Se $\theta$ é $ \psi \square \eta$, em que $\square$ é $\vee$, $\wedge$, $\rightarrow$ ou $\leftrightarrow$, então $\subst(\theta, \rho)$ é $\subst(\psi, \rho) \square \subst(\eta, \rho)$.
    \end{itemize}
\end{metadefinition}
A definição acima se justifica, pois se $\theta$ é uma fórmula proposicional cujas variáveis proposicionais que ocorrem em $\theta$ são variáveis proposicionais de $\mathcal P$, o mesmo ocorre para toda subfórmula de $\theta$.

Ainda, é imediato da definição de substituição proposicional que, se $\theta$ é uma fórmula proposicional e $\rho$ é uma atribuição de $\theta$, então $\subst(\theta, \rho)$ é uma fórmula proposicional.

\begin{metalemma}
    Seja $\mathcal F$ uma coleção finita de fórmulas proposicionais e $v$ uma valoração de cada fórmula proposicional de $\mathcal F$.

    Seja $\theta$ uma fórmula proposicional e $\rho$ uma atribuição de $\theta$ tal que, para cada variável proposicional $P$ que ocorre em $\theta$, $\rho(P)$ é uma fórmula proposicional de $\mathcal F$.

    Seja $\hat v$ a valoração de variáveis proposicionais definida de modo que para cada variável proposicional $P$ que ocorre em $\theta$, $\hat v(P)$ é $v(\rho(P))$.

    Então $v(\subst(\theta, \rho))$ é igual a $\hat v(\theta)$.
\end{metalemma}

\begin{proof}
    Agimos por indução na complexidade das fórmulas $\xi$ para mostrar que se $\xi$ é uma subfórmula de $\theta$, então $v(\subst(\xi, \rho))$ é igual a $\hat v(\xi)$.

    Suponha que a tese é verdadeira para toda subfórmula própria de $\xi$
    Veremos que a tese é verdadeira para $\xi$.
    Podemos supor, então, que $\xi$ é uma subfórmula de $\theta$.

    \textbf{Caso 1:} $\xi$ é uma variável proposicional $P$.
    Então $v(\subst(\xi, \rho))$ é $v(\rho(P))$, que é $\hat v(P)$.

    \textbf{Caso 2:} $\xi$ é $\neg \psi$.
    Por hipótese indutiva, $v(\subst(\psi, \rho))$ é igual a $\hat v(\psi)$.
    Deste modo, $v(\subst(\xi, \rho))$ é $\neg v(\subst(\psi, \rho))$, que é $\neg \hat v(\psi)$, que é $\hat v(\xi)$.

    \textbf{Caso 3:} $\xi$ é $ \psi \square \eta$, em que $\square$ é $\vee$, $\wedge$, $\rightarrow$ ou $\leftrightarrow$.
    Por hipótese indutiva, $v(\subst(\psi, \rho))$ é igual a $\hat v(\psi)$, e $v(\subst(\eta, \rho))$ é igual a $\hat v(\eta)$.
    Deste modo, $v(\subst(\xi, \rho))$ é $v(\subst(\psi, \rho) \square \subst(\eta, \rho))$, que é $v(\subst(\psi, \rho)) \square v(\subst(\eta, \rho))$, que é $\hat v(\psi) \square \hat v(\eta)$, que é $\hat v(\xi)$.
\end{proof}

\begin{metacorollary}
    Se $\theta$ é uma tautologia e $\rho$ é uma atribuição para $\theta$, então $\subst(\theta,\rho)$ é uma tautologia.
\end{metacorollary}
\begin{proof}
    Seja $v$ uma valoração de $\subst(\theta,\rho)$.
    Seja $\mathcal F=\{\rho(P): P \text{ ocorre em } \theta\}$.
    Então $v$ é uma valoração de cada fórmula proposicional de $\mathcal F$.

    Seja $\hat v$ a valoração de variáveis proposicionais definida de modo que, para cada variável proposicional $P$ que ocorre em $\theta$,
    \[
        \hat v(P)=v(\rho(P)).
    \]
    Como $\theta$ é uma tautologia, temos $\hat v(\theta)=T$.
    Pelo lema anterior,
    \[
        v(\subst(\theta,\rho))=\hat v(\theta).
    \]
    Logo, $v(\subst(\theta,\rho))=T$.
\end{proof}

\begin{metatheorem}\label{thm:PCteoremasSaoTautologias}Se $\phi$ é um teorema da lógica proposicional, então $\phi$ é uma tautologia.
\end{metatheorem}
\begin{proof}
    Primeiramente, vamos mostrar que cada axioma da lógica proposicional é uma tautologia.
    Pelo corolário anterior, é suficiente mostrar que cada axioma da lógica da Metadefinição~\ref{mdef:axiomasLogicaProposicional} é uma tautologia quando $\psi$, $\phi$ e $\chi$ são variáveis proposicionais distintas, digamos, $P$, $Q$ e $R$.

    Verificaremos isso diagramaticamente, através de tabelas-verdade.
    Na tabela, cada coluna à esquerda da última coluna corresponde a uma subfórmula do axioma, e a última coluna corresponde ao valor de verdade do axioma.
    As colunas iniciais correspondem a cada variável proposicional que ocorre no axioma, e cada linha corresponde a um caso possível de atribuição de valores de verdade a estas variáveis proposicionais.
    \begin{itemize}
        \item Para o Axioma PC1, temos:
        \begin{center}
            \begin{tabular}{c|c|c}
                $P$ & $Q$ & $P \rightarrow (Q \rightarrow P)$ \\
                \hline
                $T$ & $T$ & $T$ \\
                $T$ & $F$ & $T$ \\
                $F$ & $T$ & $T$ \\
                $F$ & $F$ & $T$
            \end{tabular}
        \end{center}
        \item Para o Axioma PC2, temos:
        \begin{center}
            \begin{tabular}{c|c|c|c|c|c}
                $P$ & $Q$ & $R$ & $P \rightarrow (Q \rightarrow R)$ & $(P\rightarrow Q) \rightarrow (P \rightarrow R)$ & (PC2) \\
                \hline
                $T$ & $T$ & $T$ & $T$ & $T$ & $T$ \\
                $T$ & $T$ & $F$ & $F$ & $F$ & $T$ \\
                $T$ & $F$ & $T$ & $T$ & $T$ & $T$ \\
                $T$ & $F$ & $F$ & $T$ & $T$ & $T$ \\
                $F$ & $T$ & $T$ & $T$ & $T$ & $T$ \\
                $F$ & $T$ & $F$ & $T$ & $T$ & $T$ \\
                $F$ & $F$ & $T$ & $T$ & $T$ & $T$ \\
                $F$ & $F$ & $F$ & $T$ & $T$ & $T$
            \end{tabular}
        \end{center}
        \item Para o Axioma PC3, temos:
        \begin{center}
            \begin{tabular}{c|c|c|c|c|c}
                $P$ & $Q$ & $P \rightarrow Q$ & $P \rightarrow \neg Q$ & $\neg P$ & (PC3) \\
                \hline
                $T$ & $T$ & $T$ & $F$ & $F$ & $T$ \\
                $T$ & $F$ & $F$ & $T$ & $F$ & $T$ \\
                $F$ & $T$ & $T$ & $T$ & $T$ & $T$ \\
                $F$ & $F$ & $T$ & $T$ & $T$ & $T$
            \end{tabular}
        \end{center}
        \item Para o Axioma PC4, temos:
        \begin{center}
            \begin{tabular}{c|c|c|c}
                $P$ & $Q$ & $(P \land Q)$ & $(P \land Q) \rightarrow P$ \\
                \hline
                $T$ & $T$ & $T$ & $T$ \\
                $T$ & $F$ & $F$ & $T$ \\
                $F$ & $T$ & $F$ & $T$ \\
                $F$ & $F$ & $F$ & $T$
            \end{tabular}
        \end{center}
        \item Para o Axioma PC5, temos:
        \begin{center}
            \begin{tabular}{c|c|c|c}
                $P$ & $Q$ & $(P \land Q)$ & $(P \land Q) \rightarrow Q$ \\
                \hline
                $T$ & $T$ & $T$ & $T$ \\
                $T$ & $F$ & $F$ & $T$ \\
                $F$ & $T$ & $F$ & $T$ \\
                $F$ & $F$ & $F$ & $T$
            \end{tabular}
        \end{center}
        \item Para o Axioma PC6, temos:
        \begin{center}
            \begin{tabular}{c|c|c|c}
                $P$ & $Q$ & $(P \land Q)$ & $P \rightarrow (Q \rightarrow (P \land Q))$ \\
                \hline
                $T$ & $T$ & $T$ & $T$ \\
                $T$ & $F$ & $F$ & $T$ \\
                $F$ & $T$ & $F$ & $T$ \\
                $F$ & $F$ & $F$ & $T$
            \end{tabular}
        \end{center}
        \item Para o Axioma PC7, temos:
        \begin{center}
            \begin{tabular}{c|c|c|c}
                $P$ & $Q$ & $(P \lor Q)$ & $P \rightarrow (P \lor Q)$ \\
                \hline
                $T$ & $T$ & $T$ & $T$ \\
                $T$ & $F$ & $T$ & $T$ \\
                $F$ & $T$ & $T$ & $T$ \\
                $F$ & $F$ & $F$ & $T$
            \end{tabular}
        \end{center}
        \item Para o Axioma PC8, temos:
        \begin{center}
            \begin{tabular}{c|c|c|c}
                $P$ & $Q$ & $(P \lor Q)$ & $Q \rightarrow (P \lor Q)$ \\
                \hline
                $T$ & $T$ & $T$ & $T$ \\
                $T$ & $F$ & $T$ & $T$ \\
                $F$ & $T$ & $T$ & $T$ \\
                $F$ & $F$ & $F$ & $T$
            \end{tabular}
        \end{center}
        \item Para o Axioma PC9, temos:
        \begin{center}
            \begin{tabular}{c|c|c|c|c|c|c}
                $P$ & $Q$ & $R$ & $P\lor Q$ & $P \rightarrow R$ & $Q \rightarrow R$ & (PC9) \\
                \hline
                $T$ & $T$ & $T$ & $T$ & $T$ & $T$ & $T$ \\
                $T$ & $T$ & $F$ & $T$ & $F$ & $F$ & $T$ \\
                $T$ & $F$ & $T$ & $T$ & $T$ & $T$ & $T$ \\
                $T$ & $F$ & $F$ & $T$ & $F$ & $T$ & $T$ \\
                $F$ & $T$ & $T$ & $T$ & $T$ & $T$ & $T$ \\
                $F$ & $T$ & $F$ & $T$ & $T$ & $F$ & $T$ \\
                $F$ & $F$ & $T$ & $F$ & $T$ & $T$ & $T$ \\
                $F$ & $F$ & $F$ & $F$ & $T$ & $T$ & $T$
            \end{tabular}
        \end{center}
        \item Para o Axioma PC10, temos:
        \begin{center}
            \begin{tabular}{c|c|c|c}
                $P$ & $Q$ & $(P \leftrightarrow Q)$ & $(P \leftrightarrow Q) \rightarrow (P \rightarrow Q)$ \\
                \hline
                $T$ & $T$ & $T$ & $T$ \\
                $T$ & $F$ & $F$ & $T$ \\
                $F$ & $T$ & $F$ & $T$ \\
                $F$ & $F$ & $T$ & $T$
            \end{tabular}
        \end{center}
        \item Para o Axioma PC11, temos:
        \begin{center}
            \begin{tabular}{c|c|c|c}
                $P$ & $Q$ & $(P \leftrightarrow Q)$ & $(P \leftrightarrow Q) \rightarrow (Q \rightarrow P)$ \\
                \hline
                $T$ & $T$ & $T$ & $T$ \\
                $T$ & $F$ & $F$ & $T$ \\
                $F$ & $T$ & $F$ & $T$ \\
                $F$ & $F$ & $T$ & $T$
            \end{tabular}
        \end{center}
        \item Para o Axioma PC12, temos:
        \begin{center}
            \begin{tabular}{c|c|c|c|c|c}
                $P$ & $Q$ & $P \rightarrow Q$ & $Q \rightarrow P$ & $P \leftrightarrow Q$ & (PC12) \\
                \hline
                $T$ & $T$ & $T$ & $T$ & $T$ & $T$ \\
                $T$ & $F$ & $F$ & $T$ & $F$ & $T$ \\
                $F$ & $T$ & $T$ & $F$ & $F$ & $T$ \\
                $F$ & $F$ & $T$ & $T$ & $T$ & $T$
            \end{tabular}
        \end{center}
        \item Para o Axioma PC13, temos:
        \begin{center}
            \begin{tabular}{c|c|c|c}
                $P$ & $\neg P$ & $\neg\neg P$ & $\neg\neg P \rightarrow P$ \\
                \hline
                $T$ & $F$ & $T$ & $T$ \\
                $F$ & $T$ & $F$ & $T$
            \end{tabular}
        \end{center}
    \end{itemize}

    Feito isso, suponha que $\phi$ é um teorema da lógica proposicional, e seja $\phi_0, \dots, \phi_n$ uma dedução formal de $\phi$ a partir da coleção vazia de fórmulas proposicionais.
    Agimos por indução para mostrar que $\phi_i$ é uma tautologia para cada $i\leq n$.
    Suponha que a hipótese indutiva é verdadeira para cada $j < i$ e que $i\leq n$.

    \textbf{Caso 1:} $\phi_i$ é um axioma da lógica proposicional.
    Então $\phi_i$ é uma tautologia, como mostrado acima.

    \textbf{Caso 2:} existem $j, k < i$ tais que $\phi_j$ é da forma $\phi_k \rightarrow \phi_i$.
    Então, pela hipótese indutiva, $\phi_k\rightarrow \phi_i$ é uma tautologia, e $\phi_k$ é uma tautologia.
    
    Considere uma valoração $v$ de $\phi_i$.
    Estenda $v$ a uma valoração de $\phi_k$, digamos, $u$.
    Deste modo, $u$ é uma valoração de $\phi_k\rightarrow \phi_i$.
    Temos que $u(\phi_k)$ é $T$, pois $\phi_k$ é uma tautologia, e $u(\phi_k\rightarrow \phi_i)$ é $T$, pois $\phi_k\rightarrow \phi_i$ é uma tautologia.

    Como $u(\phi_k)$ é $T$ e $u(\phi_k\rightarrow \phi_i)$ é $T$, então $u(\phi_i)$ é $T$ (caso contrário $u(\phi_k\rightarrow \phi_i)$ seria $F$).

    Como $u$ estende $v$, as duas valorações coincidem nas variáveis proposicionais que ocorrem em $\phi_i$.
    Logo, pelo lema anterior, $v(\phi_i)=u(\phi_i)=T$.
    Portanto, $\phi_i$ é uma tautologia.
\end{proof}

O próximo passo, mais complexo, é mostrar que toda tautologia é um teorema da lógica proposicional.

\begin{convention}\label{conv:notacaoValoracaoAssinada}
    Seja $\phi$ uma fórmula proposicional e seja $v$ uma valoração de $\phi$.

    Se $P$ é uma variável proposicional que ocorre em $\phi$, definimos:
    \[
        P^v =
        \begin{cases}
            P, & \text{se } v(P)=T,\\
            \neg P, & \text{se } v(P)=F.
        \end{cases}
    \]

    Se $\theta$ é uma subfórmula de $\phi$, definimos:
    \[
        \theta^v =
        \begin{cases}
            \theta, & \text{se } v(\theta)=T,\\
            \neg\theta, & \text{se } v(\theta)=F.
        \end{cases}
    \]

    Finalmente, se $P_0,\dots,P_{n}$ são exatamente as variáveis proposicionais que ocorrem em $\phi$, definimos:
    \[
        \Gamma_v^\phi=\{P_i^v:i\leq n\}.
    \]
    Quando a fórmula $\phi$ estiver clara pelo contexto, escreveremos apenas $\Gamma_v$.
\end{convention}

O próximo metalema formaliza a seguinte ideia: fixada uma valoração $v$, se assumirmos como hipóteses as variáveis proposicionais que são verdadeiras segundo $v$ e as negações das variáveis que são falsas segundo $v$, então conseguimos deduzir, para cada subfórmula, a própria subfórmula ou sua negação, conforme o valor que ela recebe sob $v$.

\begin{metalemma}\label{mlem:fundamentalValoracoes}
    Seja $\phi$ uma fórmula proposicional e seja $v$ uma valoração de $\phi$.
    Se $\theta$ é uma subfórmula de $\phi$, então
    \[
        \Gamma_v^\phi \vdash_P \theta^v.
    \]
\end{metalemma}

\begin{proof}
    Agimos por indução na complexidade de $\theta$.

    Fixemos a fórmula $\phi$ e a valoração $v$ de $\phi$.
    Para simplificar a notação, escreveremos $\Gamma_v$ no lugar de $\Gamma_v^\phi$.

    \textbf{Caso 1:} $\theta$ é uma variável proposicional $P$.

    Como $\theta$ é uma subfórmula de $\phi$, a variável proposicional $P$ ocorre em $\phi$.
    Pela definição de $\Gamma_v$, temos que $P^v\in \Gamma_v$.
    Como $\theta=P$, temos $\theta^v=P^v$.
    Portanto,
    \[
        \Gamma_v\vdash_P \theta^v.
    \]

    \textbf{Caso 2:} $\theta$ é $\neg\psi$.

    Pela hipótese indutiva, $\Gamma_v\vdash_P \psi^v$.
    Temos dois subcasos.

    Suponha primeiramente que $v(\psi)=F$.
    Então $v(\neg\psi)=T$.
    Logo, $\theta^v=\theta=\neg\psi$.
    Mas, como $v(\psi)=F$, temos $\psi^v=\neg\psi$.
    Pela hipótese indutiva, $\Gamma_v\vdash_P \neg\psi$.
    Portanto, $\Gamma_v\vdash_P \theta^v$.

    Agora suponha que $v(\psi)=T$.
    Então $v(\neg\psi)=F$.
    Logo,
    \[
        \theta^v=\neg\theta=\neg\neg\psi.
    \]
    Como $v(\psi)=T$, temos $\psi^v=\psi$.
    Pela hipótese indutiva,
    $\Gamma_v\vdash_P \psi$.
    Pela Metaproposição~\ref{metaprop:PC13reciproca} e modus ponens, segue que, $\Gamma_v\vdash_P \neg\neg\psi$.
    Desse modo, $\Gamma_v\vdash_P \theta^v$.

    \textbf{Caso 3:} $\theta$ é $\psi\land\eta$.
    Pela hipótese indutiva, temos $\Gamma_v\vdash_P \psi^v$ e $\Gamma_v\vdash_P \eta^v$.

    Se $v(\psi)=T$ e $v(\eta)=T$, então $v(\psi\land\eta)=T$.
    Assim, $\theta^v=\psi\land\eta$.
    Como $\psi^v=\psi$ e $\eta^v=\eta$, pela hipótese indutiva, $\Gamma_v\vdash_P \psi$ e $\Gamma_v\vdash_P \eta$.
    Por PC6, e duas aplicações sucessiva de modus ponens, temos $\Gamma_v\vdash_P \psi\land\eta$.

    Agora suponha que $v(\psi)=F$.
    Então $v(\psi\land\eta)=F$, de modo que $\theta^v=\neg(\psi\land\eta)$.
    Pela hipótese indutiva, $\Gamma_v\vdash_P \neg\psi$.

    Por PC4,
    \[
        \Gamma_v\vdash_P (\psi\land\eta)\rightarrow \psi.
    \]
    Pela monotonicidade,
    \[
        \Gamma_v\cup\{\psi\land\eta\}\vdash_P \neg\psi.
    \]
    Além disso,
    \[
        \Gamma_v\cup\{\psi\land\eta\}\vdash_P \psi,
    \]
    por PC4 e modus ponens.
    Pelo Teorema da Dedução,
    \[
        \Gamma_v\vdash_P (\psi\land\eta)\rightarrow \neg\psi.
    \]
    Finalmente, por PC3,
    \[
        \Gamma_v\vdash_P ((\psi\land \eta)\rightarrow \psi)\rightarrow (((\psi\land \eta)\rightarrow \neg\psi)\rightarrow \neg(\psi\land\eta)).
    \]
    Aplicando modus ponens duas vezes, concluímos que $\Gamma_v\vdash_P \neg(\psi\land\eta)$, e, portanto, $\Gamma_v\vdash_P \theta^v$.

    O caso em que $v(\eta)=F$ é análogo utilizando PC5 ao invés de PC4.

    \textbf{Caso 4:} $\theta$ é $\psi\lor\eta$.
    Pela hipótese indutiva, $\Gamma_v\vdash_P \psi^v$ e $\Gamma_v\vdash_P \eta^v$.

    Se $v(\psi)=T$, então $v(\psi\lor\eta)=T$.
    Logo, $\theta^v=\psi\lor\eta$.
    Como $\psi^v=\psi$, pela hipótese indutiva, $\Gamma_v\vdash_P \psi$.
    Por PC7, $\Gamma_v\vdash_P \psi\lor\eta$.
    Portanto, $\Gamma_v\vdash_P \theta^v$.

    O caso em que $v(\eta)=T$ é análogo, utilizando PC8 ao invés de PC7.

    Resta o caso em que $v(\psi)=F$ e $v(\eta)=F$.
    Então $v(\psi\lor\eta)=F$, de modo que
    \[
        \theta^v=\neg(\psi\lor\eta).
    \]
    Pela hipótese indutiva,
        $\Gamma_v\vdash_P \neg\psi$ e $\Gamma_v\vdash_P \neg\eta$.
    Mostremos que
    \[
        \Gamma_v\cup\{\psi\lor\eta\}\vdash_P \eta.
    \]
    Para tanto, por PC9, é suficiente ver que $\Gamma_v\cup\{\psi\lor\eta\}\vdash_P \psi\rightarrow \eta$ e que $\Gamma_v\cup\{\psi\lor\eta\}\vdash_P \eta\rightarrow \eta$.
    A segunda afirmação já sabemos ser verdadeira.
    Para a primeira, é suficiente ver que $\Gamma_v\cup\{\psi\lor\eta, \psi\}\vdash_P \eta$.
    Ora, pela monotonicidade, $\Gamma_v\cup\{\psi\lor\eta, \psi\}\vdash_P \neg \psi$.
    Como $\Gamma_v\cup\{\psi\lor\eta, \psi\}\vdash_P \psi$, segue do Princípio da Explosão que $\Gamma_v\cup\{\psi\lor\eta, \psi\}\vdash_P \eta$.

    Assim, $\Gamma_v\cup\{\psi\lor\eta\}\vdash_P \eta$.
    Desta forma, $\Gamma_v\vdash_P (\psi\lor\eta)\rightarrow \eta$.
    Por outro lado, por PC1 e modus ponens, $\Gamma_v\vdash_P (\psi\lor\eta)\rightarrow \neg \eta$.
    Finalmente, por duas aplicações de modus ponens e PC3, $\Gamma_v\vdash_P \neg(\psi\lor\eta)$, e, portanto, $\Gamma_v\vdash_P \theta^v$.

    \textbf{Caso 5:} $\theta$ é $\psi\rightarrow\eta$.

    Pela hipótese indutiva, $\Gamma_v\vdash_P \psi^v$ e $\Gamma_v\vdash_P \eta^v$.

    Suponha primeiramente que $v(\psi)=F$.
    Então $v(\psi\rightarrow\eta)=T$, de modo que
    \[
        \theta^v=\psi\rightarrow\eta.
    \]
    Pela hipótese indutiva, $\Gamma_v\vdash_P \neg\psi$.
    Devemos ver que $\Gamma_v\vdash_P \psi\rightarrow\eta$.
    É suficiente ver que $\Gamma_v\cup\{\psi\}\vdash_P \eta$.
    Ora, por monotonicidade, $\Gamma_v\cup\{\psi\}\vdash_P \neg\psi$, e segue que $\Gamma_v\cup\{\psi\}\vdash_P \psi$.
    Pelo Princípio da Explosão, $\Gamma_v\cup\{\psi\}\vdash_P \eta$.

    Agora suponha que $v(\eta)=T$.
    Então novamente $v(\psi\rightarrow\eta)=T$, de modo que $\theta^v=\psi\rightarrow\eta$.
    Pela hipótese indutiva, $\Gamma_v\vdash_P \eta$.
    De PC1, segue que $\Gamma_v\vdash_P \eta\rightarrow (\psi\rightarrow\eta)$.
    Por modus ponens, $\Gamma_v\vdash_P \psi\rightarrow\eta$.
    
    Resta o caso em que $v(\psi)=T$ e $v(\eta)=F$.
    Então $v(\psi\rightarrow\eta)=F$, de modo que $\theta^v=\neg(\psi\rightarrow\eta)$.
    Pela hipótese indutiva, $\Gamma_v\vdash_P \psi$
    e $\Gamma_v\vdash_P \neg\eta$.

    Temos que $\Gamma_v\vdash_P (\psi\rightarrow\eta)\rightarrow \eta$ por uma aplicação do Teorema da Dedução e modus ponens, já que $\Gamma_v\cup\{\psi\rightarrow \eta\}\vdash_P \psi$.
    Por outro lado, por PC1 e modus ponens, $\Gamma_v\vdash_P (\psi\rightarrow\eta)\rightarrow \neg\eta$.
    Assim, por PC3 e duas aplicações de modus ponens, $\Gamma_v\vdash_P \neg(\psi\rightarrow\eta)$, e, portanto, $\Gamma_v\vdash_P \theta^v$.
     
    \textbf{Caso 6:} $\theta$ é $\psi\leftrightarrow\eta$.
    Pela hipótese indutiva, $\Gamma_v\vdash_P \psi^v$ e $\Gamma_v\vdash_P \eta^v$.

    Suponha primeiramente que $v(\psi)=v(\eta)$.
    Então $v(\psi\leftrightarrow\eta)=T$, de modo que $\theta^v=\psi\leftrightarrow\eta$.

    Se $v(\psi)=v(\eta)=T$, então, pela hipótese indutiva, $\Gamma_v\vdash_P \psi$ e $\Gamma_v\vdash_P \eta$.

    Pelo Axioma PC1, $\Gamma_v\vdash_P \psi\rightarrow (\eta\rightarrow \psi)$ e $\Gamma_v\vdash_P \eta\rightarrow (\psi\rightarrow \eta)$.
    Assim, por modus ponens, $\Gamma_v\vdash_P \eta\rightarrow \psi$ e $\Gamma_v\vdash_P \psi\rightarrow \eta$.
    
    Se $v(\psi)=v(\eta)=F$, então, pela hipótese indutiva, $\Gamma_v\vdash_P \neg\psi$ e $\Gamma_v\vdash_P \neg\eta$.
    Vejamos que $\Gamma_v\vdash_P \psi\rightarrow\eta$.
    Basta ver que $\Gamma_v\cup\{\psi\}\vdash_P \eta$.
    Mas isso decorre do princípio da explosão, já que $\Gamma_v\cup\{\psi\}\vdash_P \neg\psi$ e $\Gamma_v\cup\{\psi\}\vdash_P \psi$.
    Analogamente, $\Gamma_v\vdash_P \eta\rightarrow\psi$.

        Em ambos os subcasos, temos
    \[
        \Gamma_v\vdash_P \psi\rightarrow\eta
        \qquad\text{e}\qquad
        \Gamma_v\vdash_P \eta\rightarrow\psi.
    \]
    Por PC12,
    \[
        \Gamma_v\vdash_P
        (\psi\rightarrow\eta)\rightarrow((\eta\rightarrow\psi)\leftrightarrow(\psi\leftrightarrow\eta)).
    \]
    Por modus ponens,
    \[
        \Gamma_v\vdash_P
        (\eta\rightarrow\psi)\leftrightarrow(\psi\leftrightarrow\eta).
    \]
    Por PC10,
    \[
        \Gamma_v\vdash_P
        ((\eta\rightarrow\psi)\leftrightarrow(\psi\leftrightarrow\eta))
        \rightarrow
        ((\eta\rightarrow\psi)\rightarrow(\psi\leftrightarrow\eta)).
    \]
    Aplicando modus ponens duas vezes, obtemos
    \[
        \Gamma_v\vdash_P \psi\leftrightarrow\eta.
    \]
    Portanto, $\Gamma_v\vdash_P \theta^v$.

    Agora suponha que $v(\psi)\neq v(\eta)$.
    Então $v(\psi\leftrightarrow\eta)=F$, de modo que $\theta^v=\neg(\psi\leftrightarrow\eta)$.

    Suponha primeiro que $v(\psi)=T$ e $v(\eta)=F$.
    Pela hipótese indutiva,
    $\Gamma_v\vdash_P \psi$
    e
    $\Gamma_v\vdash_P \neg\eta$.

    A partir de $\Gamma_v\cup\{\psi\leftrightarrow\eta\}$, por PC10 e modus ponens, obtemos $\Gamma_v\cup\{\psi\leftrightarrow\eta\}\vdash_P \psi\rightarrow\eta$.
    Por monotonicidade, $\Gamma_v\cup\{\psi\leftrightarrow\eta\}\vdash_P \psi$.
    Logo, por modus ponens, $\Gamma_v\cup\{\psi\leftrightarrow\eta\}\vdash_P \eta$.
    Assim, $\Gamma_v\vdash_P (\psi\leftrightarrow\eta)\rightarrow \eta$ e $\Gamma_v\vdash_P (\psi\leftrightarrow\eta)\rightarrow \neg\eta$.
    Por PC3 e duas aplicações de modus ponens, $\Gamma_v\vdash_P \neg(\psi\leftrightarrow\eta)$, e, portanto, $\Gamma_v\vdash_P \theta^v$.
    
    O caso em que $v(\psi)=F$ e $v(\eta)=T$ é análogo, utilizando PC11 ao invés de PC10.

    Isso encerra todos os casos possíveis para $\theta$.
    Portanto, por indução na complexidade de $\theta$, se $\theta$ é uma subfórmula de $\phi$, então $\Gamma_v^\phi\vdash_P \theta^v$.
\end{proof}

Para concluirmos nosso objetivo, precisamos do seguinte princípio auxiliar, que será usado para provar a quebra de casos.

\begin{metalemma}\label{mlem:terceiroExcluidoProposicional}
    Seja $\phi$ uma fórmula proposicional.
    Então $\vdash_P \phi\vee \neg \phi$.
\end{metalemma}
\begin{proof}
    Escreva $A$ para $\phi\vee\neg\phi$.

    Por PC7, temos $\vdash_P \phi\rightarrow A$.
    Pelo item (a) da Metaproposição~\ref{metaprop:modusTollensProposicional} e modus ponens,
    \[
        \vdash_P \neg A\rightarrow \neg\phi.
    \]

    Por PC8, temos $\vdash_P \neg\phi\rightarrow A$.
    Novamente pelo item (a) da Metaproposição~\ref{metaprop:modusTollensProposicional} e modus ponens,
    \[
        \vdash_P \neg A\rightarrow \neg\neg\phi.
    \]

    Como $\vdash_P \neg\neg\phi\rightarrow\phi$ por PC13, segue, pelo Teorema da Dedução e por modus ponens, que
    \[
        \vdash_P \neg A\rightarrow \phi.
    \]

    Por PC3,
    \[
        \vdash_P (\neg A\rightarrow \phi)\rightarrow((\neg A\rightarrow \neg\phi)\rightarrow\neg\neg A).
    \]
    Aplicando modus ponens duas vezes, obtemos
    \[
        \vdash_P \neg\neg A.
    \]
    Por PC13 e modus ponens,
    \[
        \vdash_P A.
    \]
    Como $A$ é $\phi\vee\neg\phi$, segue que $\vdash_P \phi\vee\neg\phi$.
\end{proof}

\begin{metalemma}[Quebra de casos]\label{mlem:quebraDeCasosProposicional}
    Seja $\Gamma$ uma coleção de fórmulas proposicionais, e sejam $\phi$ e $\psi$ fórmulas proposicionais.
    Se $\Gamma\cup\{\phi\}\vdash_P \psi$ e $\Gamma\cup\{\neg \phi\}\vdash_P \psi$, então $\Gamma\vdash_P \psi$.
\end{metalemma}
\begin{proof}
    Pelo Teorema da Dedução, basta ver que se
    \[
        \Gamma\vdash_P \phi\rightarrow\psi
        \qquad\text{e}\qquad
        \Gamma\vdash_P \neg\phi\rightarrow\psi,
    \]
    então $\Gamma\vdash_P \psi$.

    Por PC9,
    \[
        \Gamma\vdash_P
        (\phi\rightarrow\psi)\rightarrow((\neg\phi\rightarrow\psi)\rightarrow((\phi\vee\neg\phi)\rightarrow\psi)).
    \]
    Aplicando modus ponens duas vezes,
    \[
        \Gamma\vdash_P (\phi\vee\neg\phi)\rightarrow\psi.
    \]
    Pelo lema anterior e pela monotonicidade,
    \[
        \Gamma\vdash_P \phi\vee\neg\phi.
    \]
    Logo, por modus ponens,
    \[
        \Gamma\vdash_P \psi.
    \]
\end{proof}
\begin{metatheorem}[Teorema das tautologias]\label{thm:tautologiasSaoTeoremas}
    Se $\phi$ é uma tautologia, então
    \[
        \vdash_P \phi.
    \]
\end{metatheorem}
A estratégia da prova é começar com atribuições completas de valores de verdade, caso em que o Metalema~\ref{mlem:fundamentalValoracoes} se aplica diretamente, e então remover as hipóteses uma a uma por meio da quebra de casos.
\begin{proof}
    Suponha que $\phi$ é uma tautologia.
    Seja $Q_0,\dots,Q_n$
    uma enumeração das variáveis proposicionais que ocorrem em $\phi$.

    Para cada $k\leq n+1$, chamaremos de \emph{$k$-atribuição parcial} uma atribuição de valores de verdade para as variáveis $Q_0,\dots,Q_{k-1}$.
    No caso $k=0$, isto significa uma atribuição cujo domínio é vazio.

    Se $s$ é uma $k$-atribuição parcial, definimos, para cada $i<k$,
    \[
        Q_i^s=
        \begin{cases}
            Q_i, & \text{se } s(Q_i)=T,\\
            \neg Q_i, & \text{se } s(Q_i)=F.
        \end{cases}
    \]
    Definimos também $\Gamma_s$ como $\{Q_i^s:i<k\}$.
    Provaremos, por indução decrescente em $k$, a seguinte afirmação:
    \[
        (\ast_k)\qquad
        \text{se }s\text{ é uma }k\text{-atribuição parcial, então }
        \Gamma_s\vdash_P \phi.
    \]
    Começamos com o caso $k=n+1$.

    Seja $s$ uma $(n+1)$-atribuição parcial.
    Então $s$ atribui valores de verdade a todas as variáveis proposicionais que ocorrem em $\phi$.
    Portanto, $s$ é uma valoração de $\phi$.

    Como $\phi$ é uma tautologia, temos
    \[
        s(\phi)=T.
    \]
    Logo, pela definição de $\phi^s$, temos $\phi^s$ é $\phi$.
    Pelo Metalema~\ref{mlem:fundamentalValoracoes}, segue que $\Gamma_s\vdash_P \phi$.
    
    Agora suponha que $k\leq n$ e que $(\ast_{k+1})$ já foi provada.
    Mostraremos que $(\ast_k)$ é verdadeira.

    Seja $s$ uma $k$-atribuição parcial.
    Devemos provar que
    \[
        \Gamma_s\vdash_P \phi.
    \]

    Considere a extensão $s_T$ de $s$ às variáveis $Q_0,\dots,Q_k$ que satisfaz $s_T(Q_k)=T$.

    Pela hipótese indutiva $(\ast_{k+1})$, temos $\Gamma_{s_T}\vdash_P \phi$.
    Mas, pela definição de $\Gamma_{s_T}$, temos $\Gamma_{s_T}=\Gamma_s\cup\{Q_k\}$.
    Portanto,
    \[
        \Gamma_s\cup\{Q_k\}\vdash_P \phi.
    \]

    Agora considere a extensão $s_F$ de $s$ às variáveis $Q_0,\dots,Q_k$ que satisfaz $s_F(Q_k)=F$.
    Então $s_F$ também é uma $(k+1)$-atribuição parcial.
    Pela hipótese indutiva $(\ast_{k+1})$, temos $\Gamma_{s_F}\vdash_P \phi$.
    Mas, pela definição de $\Gamma_{s_F}$, $\Gamma_{s_F}=\Gamma_s\cup\{\neg Q_k\}$.
    Portanto,
    \[
        \Gamma_s\cup\{\neg Q_k\}\vdash_P \phi.
    \]

    Do Metalema~\ref{mlem:quebraDeCasosProposicional}, decorre que
    \[
        \Gamma_s\vdash_P \phi.
    \]

    Assim, $(\ast_0)$ é verdadeira.
    Existe apenas uma $0$-atribuição parcial, a atribuição vazia.
    Para essa atribuição $s$, temos $\Gamma_s=\varnothing$ por vacuidade.
    Logo, $\vdash_P \phi$.
\end{proof}
Isso conclui que as tautologias são exatamente os teoremas da lógica proposicional.

Uma consequência é a seguinte.

\begin{metadefinition}
    Seja $\Gamma$ uma coleção finita de fórmulas proposicionais, e $\psi$ uma fórmula proposicional.
    Dizemos que $\psi$ é uma consequência tautológica de $\Gamma$, e escrevemos $\Gamma\models \psi$, se, e somente se, para toda valoração $v$ de $\psi$ e de cada fórmula $\xi\in \Gamma$, se $v(\xi)=T$ para cada $\xi\in \Gamma$, então $v(\psi)=T$.
\end{metadefinition}

\begin{metacorollary}
    Seja $\Gamma$ uma coleção finita de fórmulas proposicionais, e seja $\psi$ uma fórmula proposicional.
    
    Então $\Gamma\vdash_P \psi$ se, e somente se, $\Gamma\models \psi$.
\end{metacorollary}

\begin{proof}
    Procedemos por indução na quantidade de fórmulas da coleção finita $\Gamma$.

    Para o caso em que $\Gamma$ é vazia, este resultado nos diz que $\psi$ é um teorema da lógica proposicional se, e somente se, $\psi$ é uma tautologia, o que já foi mostrado.

    Suponha que o resultado é verdadeiro para toda coleção finita de $n$ fórmulas proposicionais.
    Seja $\Gamma$ uma coleção finita de $n+1$ fórmulas proposicionais, e seja $\phi$ uma fórmula de $\Gamma$.
    Seja $\Gamma'$ a coleção obtida removendo $\phi$ de $\Gamma$.
    Assim, $\Gamma=\Gamma'\cup\{\phi\}$, e $\Gamma'$ tem $n$ fórmulas.

    Suponha primeiro que $\Gamma\vdash_P \psi$.
    Seja $v$ uma valoração de $\psi$ e de $\Gamma$ tal que $v(\xi)=T$ para cada $\xi\in \Gamma$.
    Devemos ver que $v(\psi)=T$.

    Temos que $\Gamma'\cup\{\phi\}\vdash_P \psi$.
    Pelo Teorema da Dedução, $\Gamma'\vdash_P \phi\rightarrow \psi$.
    Pela hipótese indutiva, $\Gamma'\models \phi\rightarrow \psi$, de modo que $v(\phi\rightarrow \psi)=T$.
    Como $v(\phi)=T$ (pois $\phi\in \Gamma$), segue que $v(\psi)=T$.

    Agora suponha que $\Gamma\models \psi$.
    Devemos ver que $\Gamma\vdash_P \psi$.
    Equivalentemente, veremos que $\Gamma'\vdash_P \phi\rightarrow \psi$.
    Por hipótese indutiva, basta ver que $\Gamma'\models \phi\rightarrow \psi$.
    Seja $v$ uma valoração de $\phi\rightarrow \psi$ e de cada fórmula de $\Gamma'$ tal que $v(\chi)=T$ para cada $\chi\in \Gamma'$.

    Se $v(\phi)=F$, então $v(\phi\rightarrow \psi)=T$ e termina a demonstração.

    Se $v(\phi)=T$, então $v(\chi)=T$ para cada $\chi\in \Gamma$, de modo que $v(\psi)=T$, por hipótese.
    logo, $v(\phi\rightarrow \psi)=T$, e termina a prova.
\end{proof}

\subsection*{Exercícios}

\begin{exercise}
    Decida quais das seguintes sequências de símbolos são fórmulas proposicionais na notação polonesa.
    Quando forem, escreva-as na notação usual.
    Quando não forem, explique em que ponto a formação da fórmula falha.
    \begin{enumerate}[label=(\alph*)]
        \item $\rightarrow P_0 \wedge P_1 P_2$.
        \item $\wedge \rightarrow P_0 P_1 \neg P_2$.
        \item $\leftrightarrow \vee P_0 P_1 \wedge P_2 P_3$.
        \item $\neg \rightarrow P_0 \vee P_1 P_2$.
        \item $\rightarrow \neg P_0$.
        \item $\wedge P_0 \vee P_1$.
    \end{enumerate}
\end{exercise}

\begin{exercise}
    Segundo as convenções de escrita adotadas nesta seção, diga quais das seguintes expressões possuem leitura unívoca e quais precisam de parênteses adicionais.
    Nos casos em que a leitura for unívoca, escreva a fórmula completamente parentizada.
    \begin{enumerate}[label=(\alph*)]
        \item $\neg P_0\vee P_1\rightarrow P_2$.
        \item $\neg(P_0\vee P_1)\leftrightarrow P_2$.
        \item $P_0\rightarrow P_1\rightarrow P_2$.
        \item $P_0\wedge P_1\vee P_2$.
        \item $(P_0\wedge P_1)\rightarrow P_2$.
    \end{enumerate}
\end{exercise}

\begin{exercise}
    Seja $\Gamma$ uma coleção de fórmulas proposicionais.
    Mostre que, se $\Gamma\vdash_P \phi$ e $\Gamma\vdash_P \psi$, então
    \[
        \Gamma\vdash_P \phi\land\psi.
    \]
\end{exercise}

\begin{exercise}
    Use o Teorema da Dedução para mostrar que
    \[
        \vdash_P (\phi\land\psi)\rightarrow(\psi\land\phi).
    \]
\end{exercise}

\begin{exercise}
    Mostre que
    \[
        \vdash_P ((\phi\rightarrow\psi)\land(\psi\rightarrow\chi))\rightarrow(\phi\rightarrow\chi).
    \]
\end{exercise}

\begin{exercise}
    Mostre que as seguintes fórmulas são tautologias.
    Depois, usando o Teorema das Tautologias, conclua que são teoremas da lógica proposicional.
    \begin{enumerate}[label=(\alph*)]
        \item $(\phi\land\psi)\leftrightarrow(\psi\land\phi)$.
        \item $(\phi\lor\psi)\leftrightarrow(\psi\lor\phi)$.
        \item $((\phi\land\psi)\land\chi)\leftrightarrow(\phi\land(\psi\land\chi))$.
        \item $((\phi\lor\psi)\lor\chi)\leftrightarrow(\phi\lor(\psi\lor\chi))$.
    \end{enumerate}
\end{exercise}

\begin{exercise}
    Encontre uma valoração que mostre que cada uma das seguintes fórmulas não é uma tautologia.
    \begin{enumerate}[label=(\alph*)]
        \item $\phi\rightarrow\psi$.
        \item $(\phi\rightarrow\psi)\rightarrow(\psi\rightarrow\phi)$.
        \item $(\phi\lor\psi)\rightarrow\phi$.
        \item $\phi\leftrightarrow\neg\phi$.
    \end{enumerate}
\end{exercise}

\begin{exercise}
    Considere a fórmula proposicional
    \[
        \theta=(P_0\rightarrow P_1)\rightarrow(\neg P_1\rightarrow\neg P_0).
    \]
    Defina uma atribuição $\rho$ por
    \[
        \rho(P_0)=\phi\land\psi
        \qquad\text{e}\qquad
        \rho(P_1)=\chi.
    \]
    Calcule $\subst(\theta,\rho)$.
    Sabendo que $\theta$ é uma tautologia, conclua que $\subst(\theta,\rho)$ também é uma tautologia.
\end{exercise}

\begin{exercise}
    Seja $\Gamma=\{\phi\rightarrow\psi,\phi\}$, em que $\phi$ e $\psi$ são fórmulas proposicionais.
    Mostre diretamente, pela definição de consequência tautológica, que
    \[
        \Gamma\models \psi.
    \]
    Conclua que
    \[
        \Gamma\vdash_P \psi.
    \]
\end{exercise}

\begin{exercise}
    Mostre que as seguintes fórmulas são tautologias e, portanto, teoremas da lógica proposicional:
    \begin{enumerate}[label=(\alph*)]
        \item $\neg(\phi\land\psi)\leftrightarrow(\neg\phi\lor\neg\psi)$.
        \item $\neg(\phi\lor\psi)\leftrightarrow(\neg\phi\land\neg\psi)$.
        \item $(\phi\rightarrow\psi)\leftrightarrow(\neg\phi\lor\psi)$.
    \end{enumerate}
\end{exercise}